\section{Machine learning from zero}
\label{ch:ml}

\subsection{Why learn from data at all}

A program is a set of rules written by a person. For many tasks nobody can
write the rules. Nobody can list the rules that turn a German sentence into
an English one, or that say which word comes next in a paragraph, or that
tell a cat from a dog in pixels. What we do have, in enormous quantity, is
\emph{examples}: sentences with their translations, text with the word that
actually came next, pictures with labels.

Machine learning turns examples into a function by search. You fix a family
of functions with adjustable numbers (the \emph{parameters}), you write down a
score for how wrong a candidate function is on the examples (the
\emph{loss}), and you search for the parameter values that make the score
small (\emph{optimization}). ``Training'' is that search. It is the only
known way to get a translator that handles sentences no one anticipated,
because the rules come from the data, not from a person.

The price is the subject of this whole project. First, the function is only
as good as the data: a corpus that is mostly English produces a model that
is mostly English, whatever its size. Second, you cannot read the rules off
the parameters. A trained model is billions of numbers, and nobody can tell
by inspection what it learned about Persian. That is why measurement
instruments exist. LFS is one: it does not train anything, it measures what
training produced.

\subsection{The kinds of learning}

\begin{figure}[h]\centering\includegraphics[width=0.95\linewidth]{fig_ml_types.png}
\caption{The kinds of learning, where an LLM's stages sit, and where LFS sits: outside learning, as measurement of what was learned.}\label{fig:types}\end{figure}

\begin{description}[leftmargin=3.2cm, style=nextline, itemsep=3pt]
\item[Supervised learning] Every example comes with the answer: a sentence
and its translation, a question and its correct option, an image and its
label. The model learns a map from input to answer, and the loss compares its
answer to the given one. Classical machine translation and the Belebele
scoring in this project are supervised in form.
\item[Unsupervised learning] No answers, only inputs. The goal is structure:
which examples cluster together, which directions carry the most variation
(PCA), how the data are distributed. The variance decomposition behind LFS
is unsupervised analysis of this kind, applied to a model's hidden states.
\item[Self-supervised learning] The trick that made LLMs possible: manufacture
the answer from the data itself. For text, hide the next word and ask the
model to predict it. Every position in every document becomes a labeled
example for free, so training can use trillions of tokens. LLM pretraining
is self-supervised; the resulting \emph{base model} is a next-token
predictor.
\item[Reinforcement learning] No fixed answers; the model acts, receives a
reward, and adjusts to get more reward. Used in post-training (RLHF, DPO)
to turn a base model into an assistant. The proposal's Aim 2 mentions
preference training of this kind for multilinguality.
\item[Measurement] Not learning. LFS, MEXA, probes, and the logit lens
examine what a trained model contains. The moment a measurement becomes a
training signal it becomes part of the loss, and Goodhart's law applies;
Chapter~\ref{ch:lfs} shows what happened when LFS was tried as one.
\end{description}

Everything below fits one sentence: you decide the \emph{shape} of the
function (a line, a network), the data decides its \emph{parameters}, and a
\emph{loss} says how badly any candidate does. A 70-billion-parameter
language model fits that sentence too.

\subsection{Vectors, matrices, and why a hidden state is a vector}

A vector $\mathbf{x} \in \R^D$ is a list of $D$ numbers. Two operations carry
this whole project.

\textbf{Dot product.} $\mathbf{x}\cdot\mathbf{y} = \sum_d x_d y_d$. It is
large when the vectors point the same way.

\textbf{Norm and cosine.} $\|\mathbf{x}\| = \sqrt{\mathbf{x}\cdot\mathbf{x}}$
is the length. Cosine similarity, $\cos(\mathbf{x},\mathbf{y}) =
\mathbf{x}\cdot\mathbf{y} / (\|\mathbf{x}\|\|\mathbf{y}\|)$, is the dot
product after removing length; it lies in $[-1, 1]$ and is what retrieval
metrics such as MEXA compare.

A matrix $W \in \R^{m\times n}$ turns a vector in $\R^n$ into one in $\R^m$:
$W\mathbf{x}$ is $m$ dot products, one per row. A neural network is a chain
of such multiplications with simple nonlinear functions in between. A
\emph{hidden state} is just the vector that sits between two of those
multiplications; in a language model it has $D$ coordinates (1024 for
Qwen3-0.6B, 4096 for Llama-3.1-8B) and there is one per token per layer.

\begin{exercise}{cosine similarity and a first look at anisotropy}
\begin{lstlisting}
import numpy as np
rng = np.random.default_rng(0)
def cos(a, b): return a @ b / (np.linalg.norm(a) * np.linalg.norm(b))
x, y = rng.normal(size=1024), rng.normal(size=1024)
print(cos(x, y))            # near 0: random high-dimensional vectors are nearly orthogonal
big = np.zeros(1024); big[7] = 40.0   # one huge shared coordinate ("massive activation")
print(cos(x + big, y + big))          # near 1: one coordinate makes everything look similar
\end{lstlisting}
The second print is the whole reason LFS standardizes coordinates before
measuring anything (Chapter~\ref{ch:lfs}).
\end{exercise}

\subsection{Linear regression and least squares}

Predict $y$ from $x$ with a line, $\hat y = w x + b$. The residual for one
point is $y_i - \hat y_i$. The least-squares choice of $(w, b)$ minimizes
$\sum_i (y_i - \hat y_i)^2$, the sum of squared residuals. Figure~\ref{fig:linreg}
shows the fit and the residuals. $R^2$ is the fraction of the variance of
$y$ that the line explains: $1 - \sum_i (y_i-\hat y_i)^2 / \sum_i (y_i - \bar
y)^2$.

\begin{figure}[h]\centering\includegraphics[width=0.7\linewidth]{fig_ml_linreg.png}
\caption{Least squares picks the line that makes the gray residuals as small as possible in the squared sense.}\label{fig:linreg}\end{figure}

With several inputs the same idea is $\hat{\mathbf{y}} = X\mathbf{w}$ and the
closed form is $\mathbf{w} = (X^\top X)^{-1} X^\top \mathbf{y}$. Adding a
penalty $\lambda\|\mathbf{w}\|^2$ gives \emph{ridge regression}, which keeps
weights small when inputs are many or correlated. The English-hub test in the
paper is ridge regression from one language's representations to another's.

\begin{exercise}{least squares two ways}
\begin{lstlisting}
x = np.linspace(0, 10, 50); y = 1.2 * x + 2 + rng.normal(0, 1.5, 50)
X = np.column_stack([x, np.ones_like(x)])
w_closed = np.linalg.lstsq(X, y, rcond=None)[0]
r2 = 1 - ((y - X @ w_closed)**2).sum() / ((y - y.mean())**2).sum()
print(w_closed, r2)
\end{lstlisting}
Keep this; the next exercise recovers the same $\mathbf{w}$ by gradient descent.
\end{exercise}

\subsection{Gradient descent}

For most models there is no closed form, so we descend. The gradient
$\nabla_\theta L$ is the vector of partial derivatives of the loss with
respect to each parameter; it points uphill. Gradient descent repeats
$\theta \leftarrow \theta - \eta\,\nabla_\theta L(\theta)$ with a small
\emph{learning rate} $\eta$ (Figure~\ref{fig:gd}). Too small and it crawls; too
large and it overshoots. Using a random \emph{batch} of examples per step
gives stochastic gradient descent; Adam adapts the step per parameter and is
what language models are trained with.

\begin{figure}[h]\centering\includegraphics[width=0.66\linewidth]{fig_ml_gd.png}
\caption{Each step moves against the slope; steps shrink automatically as the slope flattens.}\label{fig:gd}\end{figure}

\begin{exercise}{gradient descent for the same line}
\begin{lstlisting}
w = np.zeros(2); eta = 0.01
for step in range(5000):
    pred = X @ w
    grad = 2 * X.T @ (pred - y) / len(y)     # derivative of mean squared error
    w -= eta * grad
print(w, w_closed)                           # should agree to a few decimals
\end{lstlisting}
Change \texttt{eta} to 0.1 and to 0.0001 and watch what happens.
\end{exercise}

\subsection{Classification: sigmoid, softmax, cross-entropy}

To predict a category, turn a score (a \emph{logit}) into a probability. For
two classes the sigmoid $\sigma(z) = 1/(1+e^{-z})$ does it; for $V$ classes
the softmax $p_v = e^{z_v} / \sum_u e^{z_u}$ does (Figure~\ref{fig:softmax}).
The loss is the negative log-probability of the true class, called
cross-entropy or negative log-likelihood (NLL): $-\log p_{\text{true}}$. It is
zero when the model is certain and right, and grows without bound as the
model assigns the truth less probability.

\begin{figure}[h]\centering\includegraphics[width=0.85\linewidth]{fig_ml_logistic.png}
\caption{Left: sigmoid. Right: softmax over a tiny vocabulary. A language model does the right-hand thing over 150{,}000 tokens at every position.}\label{fig:softmax}\end{figure}

This is exactly the language-model objective: the classes are vocabulary
tokens, the input is the text so far, and the loss is the NLL of the next
token. Perplexity is $\exp(\text{mean NLL})$; a perplexity of 20 means the
model is, on average, as uncertain as a fair choice among 20 tokens.

\begin{exercise}{softmax, cross-entropy, and a gradient check}
\begin{lstlisting}
def softmax(z): z = z - z.max(); e = np.exp(z); return e / e.sum()
def nll(z, true): return -np.log(softmax(z)[true])
z = rng.normal(size=5); t = 2
# analytic gradient of NLL wrt logits is p - onehot(true)
g = softmax(z).copy(); g[t] -= 1
eps = 1e-5; num = np.array([(nll(z + eps*np.eye(5)[i], t) - nll(z - eps*np.eye(5)[i], t)) / (2*eps) for i in range(5)])
print(np.abs(g - num).max())     # ~1e-10: the analytic gradient is right
\end{lstlisting}
The subtraction of \texttt{z.max()} is numerical hygiene; it prevents overflow,
the same class of problem that made fp16 pooling fail on massive activations.
\end{exercise}

\subsection{Generalization: the only trustworthy score is on held-out data}

A model can fit its training data perfectly and be useless. Split data into
\emph{training} (fit), \emph{validation} (choose settings), and \emph{test}
(report once). Figure~\ref{fig:overfit} shows the classic picture: as
capacity grows, training error keeps falling while held-out error turns
around. Regularization (ridge penalties, early stopping, dropout) trades a
little training fit for better held-out behavior.

\begin{figure}[h]\centering\includegraphics[width=0.9\linewidth]{fig_ml_overfit.png}
\caption{Left: three polynomials on the same points. Right: training versus held-out error as capacity grows.}\label{fig:overfit}\end{figure}

Two project habits are this principle applied to research itself. The
misalignment decomposition scores each map on \emph{held-out} sentences, and it
found that the flexible nonlinear map had been overfitting to same-document
leakage. And preregistration is early stopping for the analyst: decide the
test before you see the data.

\begin{exercise}{watch overfitting happen}
Fit polynomials of degree 1 to 12 to 15 noisy points from $\sin(2\pi x)$ with
\texttt{np.polyfit}; compute training error and error on 200 fresh points;
plot both against degree. Then repeat with 150 training points and note how
the turnaround moves.
\end{exercise}

\subsection{Neural networks in one picture}

A multilayer perceptron (MLP) is $\hat{\mathbf y} = W_2\,\phi(W_1 \mathbf x +
\mathbf b_1) + \mathbf b_2$ where $\phi$ is a nonlinearity applied
elementwise (ReLU $\max(0,z)$, or the smoother GELU). Without $\phi$ the two
matrices would collapse into one; the nonlinearity is what lets the network
represent anything (Figure~\ref{fig:mlp}). Training is gradient descent on
all the weights; \emph{backpropagation} is the chain rule applied
layer by layer to get the gradient efficiently, and libraries do it for you
(autograd).

\begin{figure}[h]\centering\includegraphics[width=0.72\linewidth]{fig_ml_mlp.png}
\caption{An MLP: two matrix multiplications with a nonlinearity between them. Each transformer block contains one of these, applied to every token position independently.}\label{fig:mlp}\end{figure}

\begin{exercise}{a two-layer network on XOR, gradients by hand}
\begin{lstlisting}
X = np.array([[0,0],[0,1],[1,0],[1,1]], float); y = np.array([0,1,1,0], float)
W1 = rng.normal(0, 1, (2, 8)); b1 = np.zeros(8); W2 = rng.normal(0, 1, 8); b2 = 0.0
for step in range(5000):
    h_pre = X @ W1 + b1; h = np.maximum(0, h_pre)          # ReLU
    z = h @ W2 + b2; p = 1 / (1 + np.exp(-z))               # sigmoid output
    loss = -np.mean(y*np.log(p+1e-9) + (1-y)*np.log(1-p+1e-9))
    dz = (p - y) / 4                                        # dLoss/dz
    dW2 = h.T @ dz; db2 = dz.sum()
    dh = np.outer(dz, W2) * (h_pre > 0)                     # chain rule through ReLU
    dW1 = X.T @ dh; db1 = dh.sum(0)
    for P, g in ((W1, dW1), (W2, dW2)): P -= 0.5 * g
    b1 -= 0.5 * db1; b2 -= 0.5 * db2
print(np.round(p, 2))       # approaches [0, 1, 1, 0]
\end{lstlisting}
Remove the ReLU (\texttt{h = h\_pre}) and confirm the network can no longer
solve XOR: a linear model cannot.
\end{exercise}

\subsection{Embeddings, PCA, and effective rank}

An \emph{embedding} turns a discrete thing (a word, a token, a language) into
a vector so that geometry can carry meaning: similar things get nearby
vectors. Token embeddings are the first layer of a language model.

To look at a cloud of vectors, use principal component analysis (PCA): center
the data, compute the covariance, take its eigenvectors. The top eigenvectors
are the directions of largest variance; the eigenvalues say how much variance
each carries. The singular value decomposition (SVD) computes the same thing
directly from the data matrix. From the spectrum you get the \emph{effective
rank}: how many directions really carry variance (entropy of the normalized
spectrum, exponentiated, or the participation ratio $(\sum\lambda)^2/\sum\lambda^2$).
The paper uses these to show that a collapsed model lost dimensions even when
its LFS ratio looked fine.

\begin{exercise}{PCA and effective rank of a synthetic representation cloud}
\begin{lstlisting}
Z = rng.normal(size=(500, 64)) @ np.diag(np.r_[np.linspace(5, 1, 8), 0.1*np.ones(56)])
Zc = Z - Z.mean(0)
lam = np.linalg.eigvalsh(Zc.T @ Zc / len(Zc))[::-1]
w = lam / lam.sum()
print("participation ratio:", lam.sum()**2 / (lam**2).sum())
print("entropy rank:", np.exp(-(w[w>0] * np.log(w[w>0])).sum()))
\end{lstlisting}
Both should say ``about 8'': the cloud lives in 8 real directions plus noise.
\end{exercise}

\subsection{Vocabulary you will hear in every meeting}

\begin{description}[leftmargin=2.2cm, style=nextline, itemsep=1pt]
\item[parameters / weights] the numbers the data chooses; \emph{hyperparameters} are the ones you choose (learning rate, size).
\item[epoch, batch, step] one pass over the data; a slice of it; one gradient update.
\item[pretraining] training a base model on raw text with the next-token loss; a \emph{checkpoint} is a saved set of weights.
\item[fine-tuning] continuing training on a smaller, targeted dataset; \emph{instruction tuning} is fine-tuning on prompt--response pairs.
\item[inference] running the trained model; \emph{forward pass} is one evaluation of the network.
\item[seed] the random-number starting point; three seeds means three independent training runs.
\item[in-distribution / held-out] data like the training data / data deliberately kept away from every choice you made.
\end{description}
